\subsection{The exact ball range of the full cubic-cover eigenspace}
\label{sec:cyclic-cover-range}

\paragraph{P-F45 (the natural full-eigenspace construction has a fixed ball range).}
Let $n\ge4$ and let $F(x_1,\ldots,x_n)$ define a smooth cubic
hypersurface in $\mathbb P^{n-1}$. Put
\[
 Y_F=\{w^3+F(x_1,\ldots,x_n)=0\}\subset\mathbb P^n,
 \qquad d=\dim Y_F=n-1,
 \qquad \sigma(w,x)=(\zeta w,x),
\]
where $\zeta$ is a primitive cube root of unity. Let $V_\zeta$
be the full $\zeta$-eigenspace in the primitive middle cohomology
$H^d(Y_F,\mathbb C)_{\rm prim}$, and let $\mathcal D_n$ be
the connected full classifying domain of polarized Hodge
filtrations on $V_\zeta\oplus\overline{V_\zeta}$, with this
deck action and these Hodge numbers. Then
\[
 h^{n-1-p,p}(V_\zeta)=\binom n{3p+2-n},
 \qquad
 h^{n-1-p,p}(V_{\zeta^2})=\binom n{3p+1-n},
 \tag{P-F45a}
\]
where a binomial coefficient outside $0,\ldots,n$ is zero.
For $n\ge4$, the domain $\mathcal D_n$ is a complex ball
if and only if $n=4$ or $n=5$. In these cases its dimensions
are four and ten, respectively. The first signatures of the
Hermitian polarization on $V_\zeta$, with the convention
that the pieces with even $p$ are positive, are
\[
\begin{array}{c|c|c}
n&\text{nonzero pairs }(p,h^{n-1-p,p})& (h_+,h_-)\\ \hline
4&(1,4),(2,1)&(1,4)\\
5&(1,1),(2,10)&(10,1)\\
6&(2,15),(3,6)&(15,6)\\
7&(2,7),(3,35),(4,1)&(8,35)\\
8&(2,1),(3,56),(4,28)&(29,56).
\end{array}
\tag{P-F45b}
\]
For $n=6$, $\mathcal D_6$ is the matrix ball of type
$I_{15,6}$, of complex dimension $90$, and is not a complex
ball. For every $n\ge7$, $\mathcal D_n$ contains a nonconstant
compact holomorphic curve isomorphic to $\mathbb P^1$;
in particular it cannot be biholomorphic to a bounded domain
in any complex affine space.

\paragraph{Proof of the Hodge formula and signatures.}
The eigenspace residue formula is van Geemen--Izadi,
\emph{Half twists and the cohomology of hypersurfaces},
Lemma~2.4, printed p.~285, with the cyclic-cover convention
of Section~2.1, pp.~283--284. It says that the $i$th
eigenspace in $H^{k-p,p}(Y_k)_{\rm prim}$ is the degree
$3(p+1)-k-1-i$ part of the Jacobian algebra of the branch
cubic in $k+1$ variables. Substituting $k=n-1$ and
$i=1,2$ gives the two degrees in (P-F45a).
For a smooth cubic the $n$ quadratic partial derivatives
have no common projective zero. They form a homogeneous
regular sequence in the polynomial ring on $n$ variables:
their ideal has height $n$, and a polynomial ring is
Cohen--Macaulay. Successive quotient by a homogeneous
nonzerodivisor of degree two multiplies its Hilbert series
by $1+t$ after canceling one factor $1-t$. Thus its
Jacobian algebra has Hilbert series
\[
                      \frac{(1-t^2)^n}{(1-t)^n}
                            =(1+t)^n.
\]
This proves (P-F45a). The same coefficient count can be
read directly from the squarefree monomial basis of the
Fermat cubic, as in the cited lemma's proof.

For clarity about the eigenspace convention, the residue
volume form acquires a factor $\zeta$ under $\sigma$.
In the Jacobian algebra of $w^3+F$ the remaining powers
of $w$ are $1,w$; these therefore contribute eigenvalues
$\zeta,\zeta^2$, respectively. In particular the first
formula does not interchange the two characters silently.

Let $Q$ denote the real bilinear polarization on primitive
middle cohomology. Since $\sigma$ preserves $Q$, the two
eigenspaces pair with one another and each is isotropic
for the bilinear pairing. The form
\[
                        h(v,u)=i^d Q(v,\overline u)
                        \quad(v,u\in V_\zeta)
\]
is Hermitian: the symmetry $Q(u,v)=(-1)^dQ(v,u)$ verifies
the conjugate symmetry directly. The Hodge--Riemann
relations give orthogonality of different Hodge pieces,
with sign $(-1)^p$ on $H^{d-p,p}(V_\zeta)$ for this
normalization. An exact primary statement is Griffiths,
\emph{Periods of integrals on algebraic manifolds: summary
of main results and discussion of open problems},
equation~(9.9), printed p.~254. Its sign convention is
$(-i)^d(-1)^{d-p}Q(v,\overline v)>0$, which is exactly
$(-1)^p h(v,v)>0$. Evaluating (P-F45a) gives (P-F45b).

\paragraph{The domains with two Hodge pieces.}
When precisely two adjacent Hodge pieces are nonzero,
the full filtration is specified by a maximal positive
or negative subspace for $h$. Choose a fixed orthogonal
splitting $\mathbb C^{a+b}=P\oplus N$ with positive
dimension $a$ and negative dimension $b$. Every negative
$b$-plane projects isomorphically onto $N$: a nonzero
vector in the kernel would be both positive and negative.
It is therefore the graph of a unique map $Z:N\to P$,
and negativity is equivalent to $I-Z^*Z>0$. This gives
holomorphic coordinates when the higher Hodge piece is
negative. If the higher piece is positive, use instead
its graph $P\to N$; transposing its matrix gives the
same displayed matrix domain holomorphically. Thus in
either case the domain is
\[
              \{Z\in\operatorname{Mat}_{a\times b}(\mathbb C):
                        \|Z\|_{\rm op}<1\}.
 \tag{P-F45c}
\]
For $n=4,5$ one of $a,b$ is one, so the operator norm is
the Euclidean norm and (P-F45c) is the four- or ten-ball.
For $n=6$ the two dimensions are $15,6$; its dimension
is $15\cdot6=90$.

Here is a direct proof that this last matrix domain is
not biholomorphic to a ball, without relying on a
classification of symmetric domains. For a domain $D$
containing the origin, define the intrinsic infinitesimal
seminorm
\[
 c_D(0;v)=\sup\{|df_0(v)|:f:D\to\mathbb D
                  \text{ holomorphic},\ f(0)=0\}.
\]
For the domain (P-F45c), the Schwarz lemma on the disc
$z\mapsto z v/\|v\|_{\rm op}$ gives
$c_D(0;v)\le\|v\|_{\rm op}$. Conversely, choose unit
vectors $u,w$ with $|u^*vw|=\|v\|_{\rm op}$; the function
$Z\mapsto u^*Zw$ maps $D$ to $\mathbb D$ and attains
this derivative. Thus $c_D(0;v)=\|v\|_{\rm op}$.
The identical argument for a Euclidean ball gives its
Euclidean norm. The seminorm is preserved by the
derivative of a biholomorphism, directly from its definition.

A putative biholomorphism to a ball can be composed with
a ball automorphism to send the origin to the origin.
For completeness, when $a\ne0$ the usual map
\[
 z\longmapsto
 \frac{a-P_a z-\sqrt{1-\|a\|^2}(I-P_a)z}
      {1-\langle z,a\rangle},
 \qquad P_a z=\frac{\langle z,a\rangle}{\|a\|^2}a,
\]
interchanges $0$ and $a$; expanding the squared norm
gives
$1-\|\phi_a(z)\|^2=(1-\|a\|^2)(1-\|z\|^2)
/|1-\langle z,a\rangle|^2$, and direct substitution makes
it an involution. The case $a=0$ is immediate.
The derivative would now identify the operator norm
linearly with a Hermitian norm. This is impossible when
$a,b\ge2$: for $X=E_{11}$ and $Y=E_{22}$, the operator
norms of $X,Y,X+Y,X-Y$ are all one, so the parallelogram
identity has sides two and four. Every Hermitian norm
satisfies that identity. This proves the assertion for
$n=6$.

\paragraph{Compact curves for every higher dimension.}
For $n\ge7$, let $j_0\in\{0,1,2\}$ be congruent to
$2-n$ modulo three. The three integers
$j_0,j_0+3,j_0+6$ all lie between zero and $n$.
For $n=7,8$ this is immediate, and for $n\ge9$ their
maximum is at most eight. They correspond in (P-F45a)
to three consecutive nonzero Hodge pieces, at indices
$p_0,p_0+1,p_0+2$. The first and third have the same
definite sign for $h$ and are orthogonal. Choose one
nonzero line in each and let their sum be a definite
two-dimensional subspace $S$. Keep every other Hodge
direction fixed. For every line $\ell\subset S$, place
$\ell$ in the first of these Hodge pieces and its
$h$-orthogonal complement in $S$ in the third. The
middle piece stays fixed. This gives polarized Hodge
decompositions with the same dimensions and deck action.

Their filtrations vary holomorphically with
$\ell\in\mathbb P(S)\simeq\mathbb P^1$: at every
filtration level the moving contribution is either
$0$, the line $\ell$, or the fixed space $S$. Although
the orthogonal-complement description of the lower
summand uses conjugation, the filtration itself does
not; it is the holomorphic tautological-line construction.
The conjugate eigenspace filtration is equivalently
the annihilator filtration under $Q$, which is also
holomorphic. The curve is injective, since a filtration
level between the two chosen pieces recovers $\ell$,
and lies in the connected component of the starting
Hodge structure. It is therefore a nonconstant compact
holomorphic $\mathbb P^1$ in $\mathcal D_n$.
Every holomorphic map from $\mathbb P^1$ to a domain
in $\mathbb C^N$ is constant, by the maximum principle
applied to its coordinate functions. Hence
$\mathcal D_n$ is not biholomorphic to a bounded domain,
and in particular not to a ball. Together with the
two-piece cases this proves the exact claimed range.
\hfill$\square$

\paragraph{Scope for the theta comparison.}
The domain in this proposition is the \emph{full}
polarized eigenspace classifying domain of this particular
cyclic triple cover. It is not the horizontal period image
of a restricted family. The result explains why the actual
four- and ten-ball constructions used for cubic surfaces
and cubic threefolds do not, with the same full-eigenspace
definition, give a sequence of balls as $n$ grows.
It excludes neither special ball subdomains, selected
sub-Hodge structures, half twists, different covers, nor
other orthogonal or unitary theta constructions. It supplies
no obstruction to those alternatives and no positivity or
complexity bound. The fixed-$n=5$ identification in P-F41
is unaffected.

The primary sources used for the Hodge premises are
\begin{itemize}
\item van Geemen--Izadi, Section~2.1, pp.~283--284,
Lemma~2.4, p.~285 (the exact cyclic-cover eigenspace formula):
\url{https://mathweb.ucsd.edu/~eizadi/papers/MathZ2002.pdf}.
\item Griffiths, Section~2, pp.~234--235 (the full classifying
domain and its compact flag fibers), and equation~(9.9),
p.~254 (the polarization relations):
\url{https://publications.ias.edu/sites/default/files/periodsofintegral.pdf}.
\end{itemize}
